
\begin{landscape}
%\setlength{\textwidth}{23.0cm}


\section{How to read the diagnostics}

The diagnostics is the ASCII readable text printed on unit 6
{\it stdout} (``out1000'' in Chapter~\ref{ch:startrun}) that
gives a brief overview of the global status and progress of
the cluster simulation.
Different routines write into that file, depending on the options
chosen as the input variables. The following lines occur:\\

\noindent
\hrule
\begin{minipage}{18cm}
{\tiny
\begin{verbatim}
               N  NFIX  NCRIT  NRAND  NNBOPT  NRUN

            1000     5     10   1006      50     1


             ETAI      ETAR      RS0       DTADJ     DELTAT     TCRITp    TCRIT     QE        RBAR       ZMBAR

             1.0E-02   2.0E-02   3.0E-01   1.0E+01   1.0E+01   1.0E+06   2.0E+01   2.0E-05   1.0E+00   7.0E-01


            OPTIONS

           1  2  3  4  5  6  7  8  9 10 11 12 13 14 15 16 17 18 19 20 21 22 23 24 25 26 27 28 29 30 31 32 33 34 35 36 37 38 39 40

           1  1  1  0  1  0  4  0  0  2  1  0  0  0  1  1  1  0  0  0  1  0  1  0  0  2  0  0  0  2  0  0  2  0  1  0  1  1  0  1


            OPTIONS BK:

           1  2  3  4  5  6  7  8  9 10

           0  0  0  0  0  0  0  0  0  0


            DTMIN     RMIN      ETAU      ECLOSE    GMIN      GMAX

            1.0E-04   1.0E-02   1.0E-01   1.0E+00   1.0E-06   1.0E-02

  ****** NOTE: new random number seed initialisation!
  ****** AND new ran2 from new ed. of Press et al.
\end{verbatim}
}
\end{minipage} \hfill
%
\begin{minipage}{4.5cm}
\begin{center}
written by the routine:
{\tt input.F}\\[2ex]
Usage: Repetition of the input variables
\end{center}
\end{minipage}

\vspace{2ex}
\noindent
\hrule


\vspace{1ex}
\noindent
\begin{minipage}{18cm}
{\tiny
\begin{verbatim}
            STANDARD IMF    ALPHA = 2.35  BODY1 = 20.0  BODYN = 0.10 ZMASS = 3.36752E+02 NBIN0=    0 ZMET = 0.00 EPOCH0 = 0.00

            ...............................................................................................................

            BINARY STAR IMF:   NB =  400  RANGE =  3.27E+01  2.14E-01 ZMB = 3.74E+02  <MB> = 9.36E-01

            SINGLE STAR IMF:   NS = 1200  RANGE =  7.76E+00  1.01E-01 ZMS = 5.25E+02  <MS> = 4.38E-01
\end{verbatim}
}
%{\it These lines if routine {\tt imf2.f} is used (KZ(20)$\ge$2):}
IMF power law index, max.~mass, min.~mass, total mass,
\# of primordial bin., metallicity,
evolution.~epoch [Myrs].\\
..................................................................\\
number of objects, mass range, average mass before scaling.
\end{minipage} \hfill
%
\begin{minipage}{4.5cm}
\begin{center}
{\tt data.F},\\
(if KZ(20)=0 \& BODY1$\ne$BODYN)\\[1ex]
or\\[1ex]
{\tt imf2.F}, if KZ(20)$\ge$2\\[2ex]
Information about initial mass function (IMF).
%(only if an IMF is chosen)
\end{center}
\end{minipage}

\vspace{2ex}
\noindent
\hrule


\vspace{1ex}
\noindent
\begin{minipage}[t]{18cm}
{\tiny
\begin{verbatim}
            SCALING:   SX =  1.00421D+00  E = -2.49E-01  M(1) = 5.94E-02  M(N) = 2.97E-04  <M> = 1.00E-03

            TIME SCALES:   TRH = 2.8E+01  TCR = 2.8E+00  2<R>/<V> = 2.8E+00


     PHYSICAL SCALING:    R* = 1.0000E+00  M* = 7.0000E+02  V* = 1.7348E+00  T* = 5.6466E-01  <M> = 7.0000E-01
                          SU = 4.4335E+07  AU = 2.0627E+05  YRS = 3.5408E+06
\end{verbatim}
}
scaling factor for energy, total energy,
max.~mass, min.~mass, average mass {\it after scaling};\\
Spitzer's half-mass relaxation time, crossing time obtained from
total energy and mass, crossing time obtained from virial radius
(see \ref{ch:nbtime});\\
information about physical scaling: values of one $N$--body unit
in length (pc), mass (solar masses), velocity (km/s), time (million years),
average mass of particles (solar massses), astronomical units (one
$N$--body unit) and years (one $N$--body unit).
\end{minipage} \hfill
%
\begin{minipage}{4.5cm}
\begin{center}
{\tt scale.F, units.f}
%(scaling of the cluster properties)
\end{center}
\end{minipage}

\vspace{2ex}
\noindent
\hrule


\vspace{1ex}
\noindent
\begin{minipage}{18cm}
{\tiny
\begin{verbatim}
fpoly1 time=  0.1200000035762785     
fpoly2 time=  0.2100000062584875     
\end{verbatim}
}
CPU (wall clock in parallel execution) time for initialising
the force and its time derivative ({\tt fpoly1, fpoly\_mpi.f}) 
and the second and third time derivative of the force
({\tt fpoly2, fpoly2\_mpi.f}).
The mpi-versions are called for initialisation in case of
parallel runs.
\end{minipage} \hfill
%
\begin{minipage}{4.5cm}
\begin{center}
{\tt start.F}
\end{center}
\end{minipage}

\vspace{2ex}
\noindent
\hrule


\vspace{1ex}
\noindent
\begin{minipage}[t]{18cm}
{\tiny
\begin{verbatim}
    TIME  M/MT:    1.00D-02  2.00D-02  5.00D-02  1.00D-01  2.00D-01  3.00D-01  4.00D-01  5.00D-01  7.00D-01  9.00D-01  1.00D+00 <RC
      0.0 RLAGR:   1.52D-01  1.81D-01  1.91D-01  2.83D-01  4.33D-01  5.22D-01  6.17D-01  7.52D-01  1.16D+00  1.96D+00  5.86D+00  2.97D-01
      0.0 AVMASS:  6.26D-04  4.26D-03  6.13D-04  9.45D-04  7.82D-04  1.11D-03  1.09D-03  8.35D-04  9.19D-04  1.25D-03  9.01D-04  1.23D-03
      0.0 NPARTC:        17         9         2        53       130        90        91       122       216       163       107        90
      0.0 SIGR2:   2.19D-01  1.70D-01  7.37D-01  2.44D-01  1.89D-01  2.33D-01  2.16D-01  2.43D-01  1.72D-01  8.36D-02  5.09D-02  2.09D-01
      0.0 SIGT2:   2.01D-01  8.25D-02  4.46D-02  2.12D-01  3.39D-01  2.42D-01  1.71D-01  1.83D-01  1.51D-01  1.03D-01  6.16D-02  1.63D-01
      0.0 VROT:   -8.34D-02  4.41D-01 -6.28D-01  3.14D-02 -1.54D-01 -9.44D-02 -6.45D-02  1.17D-02  5.35D-02  1.98D-02 -3.46D-02  4.46D-01
\end{verbatim}
}
Time, specification of the Lagrangian radii, core radius\\
Time, Lagrangian radii, core radius (if primordial binaries: separately
 for singles and binaries, not shown above)\\
Time, average mass between Lagrangian radii, ~ avmass in the core\\
Time, number of particles within the shell, ~ in the core\\
Time, radial velocity dispersion within the shell, ~ in the core\\
Time, tangential vel.~dispersion within the shell, ~ in the core\\
Time, rotational vel. within the shell, ~ in the core (not shown above)\\
\end{minipage} \hfill
%
\begin{minipage}{4.5cm}
\begin{center}
{\tt lagr.F}
%(Lagrangian data)
\end{center}
\end{minipage}

\vspace{2ex}
\noindent
\hrule


\vspace{1ex}
\noindent
\begin{minipage}[t]{18cm}
{\tiny
\begin{verbatim}
  0 ADJUST:  TIME =    1.00000D+01  T[Myr] =     5.65  Q = 0.52  DE =  -1.403819E-05  E =  -2.500038E-01 EBIN=   0.000000E+00 EMERGE=   0.000000E+00
\end{verbatim}
}
rank, ``ADJUST:'', total time in NB units, physical time,
virial ratio, relative energy error, total energy,
total energy of regularized pairs, energy of mergers
\end{minipage} \hfill
%
\begin{minipage}{4.5cm}
\begin{center}
{\tt adjust.F}
\end{center}
\end{minipage}

\vspace{1ex}
\noindent
\hrule


\vspace{1ex}
\noindent
\begin{minipage}[c]{18cm}
{\tiny
\begin{verbatim}
  RMIN = 1.1E-03  DTMIN = 3.5E-05 RHOM = 3.5E+02 RSCALE = 9.5E-01 RSMIN = 2.2E-01  ECLOSE = 1.05  TC =    3
  PE  N       ttot      treg      tirr      tpredtot  tint      tinit      tks      ttcomm    tadj     tmov      tprednb  tsub     tsub2    xtsub1   xtsub2
   0  1000     41.46000     29.54      7.23      0.63     40.39      0.99      0.07      0.00      0.59      0.00      1.50      0.00      0.00  0.00000D+00  0.00000D+00
\end{verbatim}
}
close encounter distance and minimum time step (for regularization search,
updated from input parameters if KZ(16)=1), maximum density, virial radius,
minimum neighbour sphere, hard binary threshold energy, 
total run time in units of initial crossing times\\
number of processors, number of particles, total processing time,
total regular processing, total irregular processing, processing of
prediction, time spent in {\tt intgrt.F}, for initialisation, for
KS integration, for communication, for adjust and energy check,
for overhead of moving data in parallel runs, for neighbour predictions,
for MPI communication after irregular (tsub) and regular (tsub2) blocks,
number of bytes transferred respectively. From xtsub1/tsub and xtsub2/tsub2
the sustained bandwidth of MPI communication can be read off.
Note, that the determination of these quantities involves a
certain overhead by many calls of {\tt cputim.F} per block,
so for critically large production runs one may want to comment
these out (most of them in {\tt intgrt.F}).
\end{minipage} \hfill
%
\begin{minipage}{4.5cm}
\begin{center}
{\tt adjust.F}
\end{center}
\end{minipage}

\vspace{1ex}
\noindent
\hrule

\vspace{1ex}
\noindent
\begin{minipage}[c]{18cm}
{\tiny
\begin{verbatim}
  0 T =   10.0  N = 1000  <NB> = 20  KS =    0  NM = 0 MM = 0 NS =  1000 NSTEPS =    1610624       273    321696       1016  DE =  -0.140382E-04  E =        -0.250004

 NRUN =  1  M# =  1  CPU = 6.91000E-01  TRC =  0.0  DMIN = 6.6E-05 6.6E-05 1.0E+02 1.0E+02  AMIN = 1.0E+02  RMAX = 0.0E+00  RSMIN = 0.22  NEFF =   128

    <R>  RTIDE  RDENS   RC      NC   MC   RHOD   RHOM  CMAX   <Cn>  Ir/R    UN    NP    RCM    VCM         AZ     EB/E   EM/E   TCR     T6
 #1 0.95   9.5   0.21  0.08      5  0.073  159.   350.    5.  37.0  0.13     0     0    0.000  0.0000   0.006197  0.000  0.000  2.83     5

       NNPRED    NBCORR  NBFULL  NBVOID  NRCONV    NICONV  NBSMIN  NBDIS  NBDIS2  NCMDER  NBDER  NFAST  NBFAST    NBLOCK     NBPRED
 #2     20204    294307       0      98    2664      9227    1576      0       0      33      0      0       0     58132    3045868

     NKSTRY  NKSREG  NKSHYP     NKSPER  NPRECT  NKSREF  NKSMOD  NTTRY  NTRIP  NQUAD  NCHAIN  NMERG  NSTEPT  NSTEPQ  NSTEPC    NBLCKR    NBFLUX
 #3   14463      45      33          0       0       0       0      0      0      0       0      0       0       0       0     10333   1903963
\end{verbatim}
}
time, actual particle number, average neighbour number, number of KS pairs,
number of merged KS pairs, number of hierarchical subsystems, number of single
stars, step numbers (irregular, irr. c.m., regular, KS), relative energy error
since last output, total energy\\
several more lines uncommented here....
\end{minipage} \hfill
%
\begin{minipage}{4.5cm}
\begin{center}
{\tt output.F}
\end{center}
\end{minipage}

\vspace{2ex}
\noindent
\hrule



\vspace{1ex}
\noindent
\begin{minipage}[t]{18cm}
{\tiny
\begin{verbatim}
 STEP I     0    3   63   91  154  220  160  133  109   44   19    4
 STEP R     0    4   77  133  249  310  179   45    3
 Max Speedup Irr:     4 3.76D+00    8 6.82D+00   16 1.14D+01   32 1.66D+01   64 2.15D+01  128 2.49D+01  256 2.66D+01  512 2.74D+01 1024 2.77D+01
 Max Speedup Reg:     4 3.71D+00    8 6.69D+00   16 1.12D+01   32 1.67D+01   64 2.22D+01  128 2.62D+01  256 2.91D+01  512 3.04D+01 1024 3.11D+01
\end{verbatim}
}
histogram of distribution of irregular (STEP I), regular (STEP R)\\
If there are p step distribution (not appearing here, STEP U,
in physical time), statistics of parallel work for irr. and reg.
steps, figures given are theoretical speedups for infinitely
fast communication (limit of large block sizes)
\end{minipage} \hfill
%
\begin{minipage}{4.5cm}
\begin{center}
{\tt levels.f}
\end{center}
\end{minipage}

\vspace{2ex}
\noindent
\hrule



\vspace{1ex}
\noindent
\begin{minipage}[c]{18cm}
{\tiny
\begin{verbatim}
         END RUN    TIME[Myr] =   11.29  TOFF/TIME/TTOT=      0.00000000     20.00000000     20.00000000  CPUTOT =    1.6  ERRTOT =-5.15000D-05  DETOT =-1.28197D-05


        0 INTEGRATION INTERVAL =   20.00    NIRR=    3237662 NIRRB=       1245 NREG=     779010 NKS=       4625


          PER TIME UNIT: NIRR= 1.61883D+05 NIRRB= 6.22500D+01 NREG= 3.89505D+04 NKS= 2.31250D+02
  Total CPU=   97.11000289410342     
\end{verbatim}
}
This is the regular end of a run giving:
the integration time, total cumulative absolute and relative errors,
cumulative number of regular, irregular, KS steps,
the step numbers per time unit and the total CPU
(wall clock for parallel) time in minutes.\\
\end{minipage} \hfill
%
\begin{minipage}{4.5cm}
\begin{center}
\vspace{1cm}
{\tt adjust.F}
\end{center}
\end{minipage}


\noindent
\hrule
\vfill
\end{landscape}

\newpage

To check a regular stop of the run, look at the end of the
diagnostics first.
If there are failures, the line ``{\tt CALCULATION HALTED}''
appears and means that the energy conservation could not be
guaranteed.
A restart with smaller steps (ETAI, ETAR) and larger neighbour
number NNBOPT may cure the problem, but not always;
persistent problems should be reported to Rainer Spurzem.

The unix command on the output file, e.g.\\

\noindent
{\tt
  {\it homedir}> grep ADJUST out1000
}\\

\noindent
produces an overview of the accuracy (energy error at every
DTADJ interval).
It may show where problems originated; a restart from the last
ADJUST before the error with smaller output intervals is one
way to look after it.
Watch out, because sometimes errors are not reproducible,
because changes in ADJUST intervals change frequencies of
prediction and small differences can build up.
A quick possibility to see the real evolution of the system
is to {\tt grep} for the lines with Lagrangian radii and other
quantities (see above), which can directly be plotted, e.g.\
with gnuplot, because the first column is always the time.


